
\subsubsection{3.1 Information Bandwidth Framework}

Reward granularity is framed in information-theoretic terms with approximate bit estimates:

\begin{itemize}
\item \textbf{LOW (Binary):} r $\in$ {0, 1} provides approximately 1 bit per update.
\item \textbf{MEDIUM (Categorical):} Error-type scoring provides ordered feedback with approximately 1.5 bits.
\item \textbf{HIGH (Bandwidth):} Composite of categorical, pass ratio, and partial credit provides approximately 2 bits.
\end{itemize}

These bit estimates are heuristic approximations; precise calculations would require entropy analysis of actual reward distributions.

\subsubsection{3.2 Error-Type Scoring}

The categorical reward assigns scores based on error type, reflecting distance to correct solution:

\begin{table}[htbp]
\centering
\resizebox{\columnwidth}{!}{%
\begin{tabular}{lll}
\toprule
Error Type & Score & Rationale \\
\midrule
Syntax error & 0.00 & Code does not parse \\
Runtime error & 0.25 & Code runs but crashes \\
Assertion failure & 0.50 & Code runs but produces wrong output \\
Pass & 1.00 & Code is correct \\
\bottomrule
\end{tabular}
}
\end{table}

\subsubsection{3.3 High-Bandwidth Reward Composition}

The high-bandwidth condition combines three signals with fixed weights:

r = 0.5 $\times$ categorical\_score + 0.3 $\times$ pass\_ratio + 0.2 $\times$ partial\_credit

where pass\_ratio is the fraction of test cases passed and partial\_credit measures numeric closeness for assertion failures when expected and actual values are extractable.

\subsubsection{3.4 Training Configuration}

\begin{itemize}
\item \textbf{Model:} CodeLlama-7B-Instruct (codellama/CodeLlama-7b-Instruct-hf)
\item \textbf{Training data:} MBPP sanitized split (approximately 374 problems)
\item \textbf{Evaluation:} MBPP validation subset (50 problems)
\item \textbf{Algorithm:} REINFORCE with advantage baseline (not PPO despite initial design)
\item \textbf{Epochs:} 1 (proof-of-concept scale)
\item \textbf{Seeds:} 3 per condition (9 total runs)
\item \textbf{Batch size:} 4
\item \textbf{Gradient steps:} Approximately 94 per epoch
\item \textbf{Learning rate:} 1e-5
\item \textbf{Test timeout:} 5.0 seconds
\end{itemize}
